\documentclass[a4paper,fleqn,doubleblind]{cas-dc}
\usepackage[T1]{fontenc}
\usepackage[utf8]{inputenc}
\usepackage{charter}
\renewcommand{\sfdefault}{bch}
\usepackage{textcomp}
\usepackage{amsmath,amssymb}
\usepackage[numbers,sort&compress]{natbib}
\usepackage{graphicx,booktabs,array,enumitem}
\usepackage{xurl}
\urlstyle{same}
\usepackage[switch]{lineno}
\usepackage[expansion=false]{microtype}
\def\tsc#1{\csdef{#1}{\textsc{\lowercase{#1}}\xspace}}
\renewcommand{\linenumberfont}{\normalfont\tiny\color{gray}}
\setlength{\bibsep}{1pt}
\renewcommand{\figurename}{Fig.}
\begin{document}
\let\WriteBookmarks\relax
\let\printorcid\relax
\def\floatpagepagefraction{1}
\def\textpagefraction{.001}
\shorttitle{Self-timed end-of-life alarms: timing, cost and limits}
\shortauthors{Anonymized manuscript}
\title [mode = title]{Self-timed end-of-life alarms: timing, maintenance cost and limits}
\begin{abstract}
Remaining-useful-life models are usually scored at evaluator-chosen times, whereas a replacement is triggered once. We evaluate first-warning policies that raise one alarm per unit, accepted if the true remaining life lies between 2 cycles below and 13 cycles above the reported estimate, and we separate prediction accuracy, alarm timing and maintenance cost. A gradient-boosted regressor with a tuned trigger placed 91.7\% of PHM08 alarms and 90.2\% of 927 simulated engine alarms in this window at a 20-cycle warning level. Fixed-time models reached at most 33.5\% at half of life, but a calibrated fixed-time model reached 81.2\% at 90\% of life: the contrast mainly reflects horizon and calibration. A constant estimate at the same alarm times scored 90.8\%, so the window rewards alarm timing. In a renewal--reward model, window-tuned alarms were cheaper than age replacement except on N-CMAPSS, but alarms tuned directly to cost were cheaper on some datasets and dearer on others. Retrained under protocols hashed in advance, and after correcting a time-scale leak, the recipe reached 68.7\% on N-CMAPSS, 0 of 17 on bearings, and 95.9\% on battery cells within batches but 65.0\% with a batch held out, where capacity-trend extrapolation reached 99.2\%.
\end{abstract}
\begin{keywords}
Remaining useful life \sep
Predictive maintenance \sep
Alarm timing \sep
Renewal--reward model \sep
Model evaluation \sep
Turbofan engine \sep
Lithium-ion battery
\end{keywords}
\maketitle
\linenumbers
\section{Introduction}\label{sec:1}

Data-driven prognostics estimates the remaining useful life (RUL) of a component from its condition-monitoring history \cite{lei2018,si2011}. On the standard turbofan benchmarks, estimates are usually scored at evaluator-chosen times: the last cycle of truncated test trajectories, every cycle, or fixed fractions of life \cite{ramasso2014}. A replacement, however, is triggered once. Four tasks are therefore easily conflated: \textbf{long-range prediction}, the accuracy of an estimate made far from failure; \textbf{near-failure estimation}, its accuracy when little life is left; \textbf{alarm-time selection}, the choice of the single moment at which to act; and \textbf{maintenance-policy performance}, the long-run cost of the replacements and failures that the alarms produce. A good score on one does not imply a good score on another.

This study began with a requirement of the third kind. An operator would accept a remaining-life estimate only if the true remaining life lay between 2 cycles below and 13 cycles above it, that is, if the error, predicted minus true cycles left, lay in [\textminus{}13, +2]. At half of life none of eleven fixed-time models placed more than 33.5\% of PHM08 predictions in this window. That figure describes long-range prediction under a narrow tolerance, not the models' quality near failure: the fixed-time gradient-boosted model among them, scored at 90\% of life with a calibration shift learned on separate engines, placed 81.2\% in the window (Section 6.1). What remains open is whether a policy that chooses its own moment, with the information available at that moment, can meet the window reliably, and whether doing so leads to good maintenance decisions.

We therefore study \textbf{first-warning policies}. A policy observes each unit cycle by cycle, stays silent while the unit looks healthy, and raises one alarm with a remaining-life estimate when it judges the end to be near. The alarm is scored when it is raised, and a unit that fails without an alarm counts as late. The idea is not new: alarm-based scheduling triggers maintenance when probabilistic RUL prognostics indicate imminent failure \cite{depater2022}, and decision-oriented metrics price a prognostic algorithm by the long-run cost of the maintenance decisions it induces \cite{kamariotis2024}. The question this paper addresses is narrower and empirical: \textbf{do first-warning policies tuned to an asymmetric tolerance window produce useful maintenance decisions, under which conditions, and with what limitations?}

To answer it we use a deliberately simple recipe built from standard parts, a gradient-boosted regressor with a capped target and a smoothed, shifted and debounced threshold trigger, all selected by nested cross-validation. We do not claim novelty for any of these parts. The contributions are the following.

\begin{itemize}[leftmargin=1.2em,itemsep=2pt]
\item[(i)] \textbf{An evaluation that separates the four tasks.} On PHM08 we decompose the contrast between fixed-time and self-timed evaluation into horizon, calibration and alarm choice. Two timing controls show that the window score mainly rewards \textit{when} the alarm is raised: a constant estimate at the same alarm times scored within a few percentage points of the model's own estimate on every dataset, and a health-indicator threshold alarm scored much lower.
\item[(ii)] \textbf{A cost comparison under identical information.} In a renewal--reward model with lead times and failure costs, alarms tuned to the window (with only the warning level chosen by cost) are compared per dataset with alarms tuned directly to the cost, including regularised and quantile-based variants, a health-indicator alarm and optimal age replacement.
\item[(iii)] \textbf{Evidence on generalisation, including failures.} The same recipe was retrained on three further datasets whose analysis protocols had been written and hashed in advance: simulated N-CMAPSS engines, laboratory lithium-ion cells and rolling bearings. An audit for this revision found that the time scale of the battery and bearing analyses depended on test-unit lives. We report the corrected re-analysis, with the battery results under a held-out test batch given equal weight to those within batches.
\end{itemize}

Three things are not claimed. A frozen model is never transferred between asset types: the recipe is retrained on each dataset. No field data are used. And no PHM08 leaderboard comparison is possible, because the task and the data split differ. Section 2 reviews related work; Section 3 defines the task and the cost model; Section 4 describes the recipe; Section 5 gives the experimental design and the status of each analysis; Section 6 reports results; Section 7 discusses them; Section 8 outlines future work; and Section 9 concludes.

\section{Related work}\label{sec:2}

\subsection{Remaining-life prediction and its evaluation}

The C-MAPSS simulations and the PHM08 challenge data \cite{saxena2008} are the standard engine benchmarks. The three winning PHM08 approaches were a similarity-based method \cite{wang2008}, a recurrent network \cite{heimes2008} and an ensemble of neural networks filtered by a Kalman filter \cite{peel2008}. Later work moved to long short-term memory \cite{zheng2017} and convolutional networks \cite{li2018}, usually with a remaining-life target capped near 125 cycles. N-CMAPSS adds realistic flight profiles \cite{ariaschao2021}. On tabular inputs, tree ensembles remain strong \cite{grinsztajn2022,friedman2001}. Severson et al. \cite{severson2019} cycled commercial lithium-ion cells to failure, and the PRONOSTIA platform produced run-to-failure vibration records of rolling bearings \cite{nectoux2012}. Ramasso and Saxena \cite{ramasso2014} showed how strongly reported scores depend on the choice of cut-off points and metrics. Saxena et al. \cite{saxena2010} introduced metrics that follow a prediction over time, such as the prognostic horizon and $\alpha$--$\lambda$ accuracy, and asymmetric losses that penalise late predictions are common \cite{rengasamy2020}.

\subsection{From prognostics to maintenance decisions}

Age replacement, which uses no condition information, is the classical benchmark for any condition-based policy \cite{barlow1960,dejonge2020}. Degradation models with Bayesian updating yield residual-life distributions \cite{gebraeel2005}, and dynamic predictive-maintenance frameworks act on predicted failure probabilities \cite{nguyen2019,mitici2023} or learn policies from probabilistic prognostics by reinforcement learning \cite{lee2023}. Two lines of work are closest to ours. de Pater et al. \cite{depater2022} trigger alarms from the evolution of probabilistic RUL prognostics of convolutional networks, with a safety factor, and feed the alarms into fleet maintenance scheduling for aircraft engines. Kamariotis et al. \cite{kamariotis2024} propose a decision-oriented metric, the long-run expected maintenance cost per unit time estimated from run-to-failure experiments, and use it both to compare prognostic algorithms and to optimise the parameters of heuristic maintenance policies; tuning an alarm policy to cost on run-to-failure data is therefore established. Xie et al. \cite{xie2025}, in a preprint, train the prognostic model itself against downstream maintenance cost on turbofan data.

\subsection{Position of this study}

We propose neither a new alarm rule nor a new cost metric. Our contribution is an analysis of a one-alarm-per-unit acceptance criterion with an asymmetric window: what it rewards, whether policies tuned to it give cost-effective alarms compared with policies tuned directly to the cost in the manner of Kamariotis et al. \cite{kamariotis2024}, and how the same recipe behaves when retrained on new data. Because such analyses are easily biased by model selection on evaluation data \cite{varma2006} and by leakage \cite{kapoor2023}, every setting is selected on training units only, and three of the datasets were analysed under protocols fixed in advance \cite{nosek2018}, with the limitations stated in Section 5.3.

\section{Problem formulation}\label{sec:3}

\subsection{First-warning task}

A unit \textit{i} is observed at cycles \textit{t} = 1, \ldots{}, \textit{L}\textsubscript{i}, where \textit{L}\textsubscript{i} is its life, unknown in service. The true remaining life at cycle \textit{t} is \textit{r}\textsubscript{t} = \textit{L}\textsubscript{i} \textminus{} \textit{t}. A warning policy maps the history up to \textit{t} to either \textit{silent} or \textit{warn}, together with an estimate $\hat{y}$\textsubscript{t}. The first-warning time is

\begin{equation}
t^{*} = \min\{\, t < L_i : \text{the policy warns at } t \,\},
\end{equation}

and the error of the warning is \textit{e} = $\hat{y}$\textsubscript{t*} \textminus{} \textit{r}\textsubscript{t*}. A warning is \textbf{correct} if \textminus{}13 $\le$ \textit{e} $\le$ +2, that is, exactly when

\begin{equation}
\hat{y}_{t^*} - 2 \le r_{t^*} \le \hat{y}_{t^*} + 13.
\end{equation}

It is \textbf{too early} if \textit{e} < \textminus{}13 and \textbf{too late} if \textit{e} > +2; a unit that fails without a warning is \textbf{missed} and counted with the late outcomes. The primary metric is the share of units whose first warning is correct (the in-window rate). We also report the late-or-missed and too-early rates, the true remaining life at the alarm, and the share of alarms raised with \textit{r}\textsubscript{t*} $\ge$ 13. Only those alarms can be too early, unless the estimate is negative, so a policy that warned only in the last few cycles would meet the window trivially. The window was fixed by the operator's requirement before this study; its early limit coincides with the early constant of the PHM08 score \cite{saxena2008}, and its late limit is much stricter than that score's late constant of 10.

The warning level \textit{H} expresses how much notice is wanted: the policy warns when it believes that about \textit{H} cycles are left. \textit{H} = 20 is the primary level, designated on PHM08 as the smallest of \{10, 15, 20, 30, 40\} at which more than 90\% of alarms can be too early (Section 5.3). On the battery and bearing data a cycle of the model is a block of consecutive charge--discharge cycles or vibration snapshots (Section 5.2), and we write ``unit'' for an engine, cell or bearing.

\subsection{Four tasks, and when evaluation leaks}

We separate (1) \textbf{long-range prediction}, the accuracy of $\hat{y}$\textsubscript{t} at times chosen independently of the policy, for example at a fixed fraction of life; (2) \textbf{near-failure estimation}, the accuracy of $\hat{y}$\textsubscript{t} when \textit{r}\textsubscript{t} is small, for example at the cycle where \textit{r}\textsubscript{t} = \textit{H}; (3) \textbf{alarm-time selection}, the distribution of \textit{r}\textsubscript{t*} produced by the policy; and (4) \textbf{maintenance-policy performance}, the cost rate defined in Section 3.3. A fixed fraction of life is not known in service, but evaluating a model retrospectively at such a time is not leakage: it becomes leakage only if the fraction, the life, or a rule derived from them is given to the model as an input or used to train or select it. The fixed-time reference models of Section 6.1 received none of these. Conversely, the first-warning task is not easier merely because the policy chooses its time: the choice must be made from the history alone, and a late or missed alarm is a failure.

\subsection{Maintenance cost}

Alarms are priced in a renewal--reward model \cite{barlow1960}. An alarm at cycle \textit{t}\textsuperscript{*} schedules a planned replacement $\ell$ cycles later. If \textit{r}\textsubscript{t*} > $\ell$ the replacement happens at \textit{t}\textsuperscript{*} + $\ell$ and costs 1; otherwise, or if no alarm is raised, the unit fails at \textit{L} and the failure costs $\rho$. An alarm with exactly $\ell$ cycles left counts as a failure. Units are replaced by new ones, so the long-run cost per cycle is estimated by the total cost divided by the total cycles of use over the test units. We use lead times $\ell$ $\in$ \{0, 5, 10, 15\} cycles (batteries: 0, 15, 30 and 45 cycles) and cost ratios $\rho$ $\in$ \{5, 10, 20, 50\}. These ratios are illustrative, not measured. Optimal age replacement replaces every unit at the age \textit{T} that minimises the cost rate on the training lives; it needs no lead time because the replacement age is known in advance. Perfect foresight (replacement at \textit{L} \textminus{} 1, cost 1) and run-to-failure bound the range, and we report every cost rate relative to perfect foresight.

\section{The self-timed warning recipe}\label{sec:4}

Fig. 1 shows the recipe. Its online part runs at every cycle of a unit in service; its offline part fixes four settings $\theta$ = (\textit{c}, \textit{m}, \textit{k}, \textit{s}) on training units only.

\begin{figure*}[!tbp]\centering
\includegraphics[width=\textwidth]{img/Fig1_STW_architecture.png}
\caption{The self-timed warning (STW) recipe. Online (top): regime-normalised sensors are encoded causally and mapped to an estimate $\hat{y}$\textsubscript{t}, which is smoothed and shifted to $\tilde{y}$\textsubscript{t}; an alarm is raised the first time $\tilde{y}$ stays at or below \textit{H} for \textit{k} cycles. Offline (bottom): $\theta$ = (\textit{c}, \textit{m}, \textit{k}, \textit{s}) is chosen on inner out-of-fold predictions of the training units and frozen; test units are scored once. Inset: the scoring window for \textit{e} = $\tilde{y}$\textsubscript{t*} \textminus{} \textit{r}\textsubscript{t*}.}
\end{figure*}

\textbf{Normalisation and features.} Each sensor is z-scored within its operating regime with statistics of the training units, and sensors that are constant within any regime are dropped (15 sensors remain on PHM08, FD001 and FD002, 16 on FD003 and FD004; N-CMAPSS uses the six temperature and shaft-speed channels fixed in its protocol). At cycle \textit{t} the encoder uses cycles 1, \ldots{}, \textit{t} only: an early-life baseline (the mean of the first min(20, \textit{t}) cycles), exponentially weighted means with spans 10 and 30 of the baseline-removed signals, least-squares slopes over the last 15 and 40 cycles, and the age \textit{t}. A check that deleting all data after \textit{t} leaves the prediction at \textit{t} unchanged is in Supplementary Section S4.

\textbf{Regressor.} A histogram gradient-boosted regressor \cite{pedregosa2011} with absolute-error loss (500 iterations, learning rate 0.05, 31 leaves, at least 40 samples per leaf, seed 0) is trained on all cycles of the training units with the target

\begin{equation}
y_t = \min( r_t, c ), \qquad c \in \{40, 60, 90\}.
\end{equation}

Capping the target is standard \cite{heimes2008}; a low cap concentrates the loss on the last \textit{c} cycles, where the alarm is raised.

\textbf{Trigger.} The estimate is smoothed causally and shifted, $\tilde{y}$\textsubscript{t} = EWM\textsubscript{m}($\hat{y}$)\textsubscript{t} + \textit{s}, with span \textit{m} $\in$ \{1, 3, 6\} (\textit{m} = 1: no smoothing) and \textit{s} $\in$ \{\textminus{}12, \textminus{}11.5, \ldots{}, +8\}, and the alarm is raised at

\begin{equation}
t^{*} = \min\{\, t : \tilde{y}_{\tau} \le H \ \text{for all}\ \tau = t-k+1, \ldots, t \,\},
\end{equation}

with debounce \textit{k} $\in$ \{1, 2, 3, 5\}; the reported estimate is $\tilde{y}$\textsubscript{t*}.

\textbf{Nested selection.} The training units of each outer fold are split into five inner folds. Regressors fitted on four inner folds, which also supply the normalisation statistics, predict every cycle of the fifth, giving out-of-fold predictions for all training units. For each \textit{H}, the grid (caps with \textit{c} $\ge$ \textit{H} + 20, three spans, four debounce values and 41 shifts) is searched for the setting with the most correct out-of-fold first warnings; ties go to the smaller |\textit{s}|. The regressor is then refitted on all training units, $\theta$ is frozen, and each test unit is processed once. Supplementary Table S15 lists every setting.

\subsection{The role of the shift}

Let $\varepsilon$ = EWM\textsubscript{m}($\hat{y}$)\textsubscript{t*} \textminus{} \textit{r}\textsubscript{t*} be the unshifted error at the alarm, so that \textit{e} = $\varepsilon$ + \textit{s}, and write the window as [\textit{a}, \textit{b}] = [\textminus{}13, +2], of width \textit{w} = 15. The following elementary result explains why a negative shift helps. It is stated for a \textbf{fixed} error distribution.

\textbf{Proposition 1.} Let $\varepsilon$ have a continuous distribution \textit{F} that does not depend on \textit{s} (assumption A1). Then the probability that \textit{e} = $\varepsilon$ + \textit{s} lies in [\textit{a}, \textit{b}] is

\begin{equation}
P(s) = F(b - s) - F(a - s).
\end{equation}

(i) \textit{P}(\textit{s}) $\le$ \textit{C}\textsubscript{w}(\textit{F}) = sup\textsubscript{x} [\textit{F}(\textit{x} + \textit{w}) \textminus{} \textit{F}(\textit{x})] for every \textit{s}. (ii) If \textit{F} has a density that is unimodal and symmetric about its median \textit{m}, \textit{P}(\textit{s}) is maximised at

\begin{equation}
s^{*} = (a + b)/2 - m = -5.5 - m,
\end{equation}

with maximum Pr(|$\varepsilon$ \textminus{} \textit{m}| $\le$ 7.5). The proof is in Supplementary Section S2.

The result implies that an unbiased estimator should be shifted by about \textminus{}5.5 cycles, because the window is centred at \textminus{}5.5. It is \textbf{not} a guarantee for the sequential policy: changing \textit{s} changes the stopping time \textit{t}\textsuperscript{*} and therefore the distribution of $\varepsilon$, so \textit{F} depends on \textit{s}. Section 6.5 measures how far the alarm moves when the shift changes and how far the realised rate departs from \textit{P}(\textit{s}).

\section{Experimental design}\label{sec:5}

\subsection{Data}

\begin{table*}[!tbp]
\footnotesize\setlength{\tabcolsep}{4pt}
\caption{Datasets. Life: cycles (battery: charge--discharge cycles; bearings: vibration snapshots) to failure of the run-to-failure units. $\kappa$: block length, computed from the training units of each outer split (Section 5.2).}
\begin{tabular}{@{}>{\raggedright\arraybackslash}p{\dimexpr0.1717\linewidth-2\tabcolsep\relax}>{\raggedright\arraybackslash}p{\dimexpr0.0997\linewidth-2\tabcolsep\relax}>{\raggedright\arraybackslash}p{\dimexpr0.1053\linewidth-2\tabcolsep\relax}>{\raggedright\arraybackslash}p{\dimexpr0.1662\linewidth-2\tabcolsep\relax}>{\raggedright\arraybackslash}p{\dimexpr0.1053\linewidth-2\tabcolsep\relax}>{\raggedright\arraybackslash}p{\dimexpr0.1717\linewidth-2\tabcolsep\relax}>{\raggedright\arraybackslash}p{\dimexpr0.1801\linewidth-2\tabcolsep\relax}@{}}
\toprule
Dataset & Units & Rows & Life, min--max (median) & Regimes & Time unit of the model & Status (Section 5.3) \\
\midrule
PHM08 & 218 & 45,918 & 128--357 (209.5) & 6 & cycle & development \\
FD001 & 100 & 20,631 & 128--362 (199) & 1 & cycle & additional evaluation \\
FD002 & 260 & 53,759 & 128--378 (199) & 6 & cycle & additional evaluation \\
FD003 & 100 & 24,720 & 145--525 (220.5) & 1 & cycle & additional evaluation \\
FD004 & 249 & 61,249 & 128--543 (234) & 6 & cycle & additional evaluation \\
N-CMAPSS & 99 & 7,473 & 48--100 (75) & 3 (flight class) & flight cycle & pre-specified \\
Battery \cite{severson2019} & 123 cells & 95,276 & 148--1935 (731) & 1 & block of $\kappa$ $\in$ \{3, 4\} cycles & pre-specified; $\kappa$ corrected \\
Bearings \cite{nectoux2012} & 17 bearings & 24,889 & 230--2803 (1428) & 3 (load/speed) & block of $\kappa$ $\in$ \{4, 6, 7, 9\} snapshots & pre-specified; $\kappa$ corrected \\
\bottomrule
\end{tabular}\par\vspace{3pt}{\raggedright\scriptsize C-MAPSS and PHM08: training sets only (complete trajectories); the truncated official test sets are not used. N-CMAPSS: one row per engine and flight (mean of the 1 Hz signals), 99 engines of nine subsets. Battery: 123 of 140 cell records after exclusions fixed before any model was run (split tests and censored cells; Supplementary Section S5); life = first cycle below 80\% of the nominal 1.1 Ah. Bearings: one row per 0.1-s snapshot recorded every 10 s; eight features per snapshot.\par}
\end{table*}

We use the 218 run-to-failure engines of PHM08, the training sets of C-MAPSS FD001--FD004 \cite{saxena2008}, and N-CMAPSS reduced to one row per engine and flight \cite{ariaschao2021}. The battery data are commercial LFP/graphite cells (A123 APR18650M1A, 1.1 Ah) cycled to failure under 66 fast-charging protocols in three batches that also differ in their rest periods \cite{severson2019}; the model sees seven per-cycle quantities (discharge and charge capacity, internal resistance, average, minimum and maximum temperature, charge time), z-scored as one regime as fixed in the protocol. The bearing data are the 17 complete records of the PRONOSTIA platform under three load and speed conditions \cite{nectoux2012}, reduced to root mean square, peak, kurtosis and crest factor of the horizontal and vertical acceleration of each snapshot.

\subsection{Time scale of the battery and bearing data}

Battery and bearing lives are several times longer, in their native units, than engine lives, and the grids of Section 4 were designed for engine lives. The protocols therefore aggregated consecutive cycles or snapshots into blocks of $\kappa$ = max(1, round(median life / 209.5)), where 209.5 cycles is the median PHM08 life, and applied the recipe, the warning level and the window in blocks. In cycles, the window is then [\textminus{}13$\kappa$, +2$\kappa$] and \textit{H} = 20 corresponds to 20$\kappa$ cycles.

\textbf{Audit finding.} The submitted analyses computed the median over all units of a dataset, including those later used as test units. Held-out lives could therefore influence the aggregation, the warning level and the window in physical units. For the batteries the effect is not hypothetical: the all-cell median, 731 cycles, lies just below the boundary between $\kappa$ = 3 and $\kappa$ = 4 (733.25 cycles). No other preprocessing step used test units (Supplementary Section S4). The block construction also filled missing per-cycle values backwards in time within a cell. This affected only the first cycle of 36 cells and never crossed a block boundary for $\kappa$ $\ge$ 2, but it has been replaced by forward filling, with a missing first block filled by the training-cell median.

\textbf{Correction.} $\kappa$ is now computed from the training units of each outer split, and inner folds inherit the outer $\kappa$. That $\kappa$ uses the lives of all training units, including those that serve as inner validation units, so it can still influence the selection of settings, but not the outer test estimate. This gave $\kappa$ = 4, 3, 4, 3, 3 in the five battery folds, 3, 4, 3 when batch 1, 2 or 3 was held out, and $\kappa$ = 7, 4, 9, 6, 7 in the five bearing folds (4 for the official split). Every affected analysis was rerun: battery cross-validation, leave-one-batch-out and the 119-cell analysis, the battery timing controls and costs, and both bearing analyses. Because the task in physical units now varies slightly between folds (for the batteries the window is [\textminus{}39, +6] cycles with $\kappa$ = 3 and [\textminus{}52, +8] cycles with $\kappa$ = 4), a \textbf{sensitivity analysis} holds the physical task fixed and varies only $\kappa$. For the batteries it uses \textit{H} = 60 cycles and the window [\textminus{}39, +6] cycles, with $\kappa$ $\in$ \{1, \ldots{}, 6\}; for the bearings it uses \textit{H} = 140 snapshots (about 23 min) and [\textminus{}91, +14] snapshots, with $\kappa$ $\in$ \{3, 5, 7, 10, 14\}. Caps and shift grids are converted to the same physical values.

\textbf{How an operator would choose $\kappa$.} In service, lives are unknown, but three sources of information are not: the lives of units previously run to failure (the training fleet, as here), design or warranty lives, and the notice that a replacement needs. We regard the task as defined in physical units, a notice of \textit{H} cycles and a tolerance window in cycles, and $\kappa$ as a temporal resolution. $\kappa$ can be fixed in advance from historical lives or selected by inner cross-validation. The sensitivity analysis shows how much the result depends on that choice.

\subsection{Development, pre-specified and exploratory analyses}

\begin{table*}[!tbp]
\footnotesize\setlength{\tabcolsep}{4pt}
\caption{Status of every analysis. Pre-specified: protocol written and hashed (SHA-256) before the model or a competitor was run on the data; the hashes are held by the authors and were not deposited with a registry.}
\begin{tabular}{@{}>{\raggedright\arraybackslash}p{\dimexpr0.2216\linewidth-2\tabcolsep\relax}>{\raggedright\arraybackslash}p{\dimexpr0.4542\linewidth-2\tabcolsep\relax}>{\raggedright\arraybackslash}p{\dimexpr0.3242\linewidth-2\tabcolsep\relax}@{}}
\toprule
Status & Analyses & What the status means \\
\midrule
Development & PHM08 (model, grids, warning levels and primary level \textit{H} = 20) & Choices made while looking at PHM08 results; estimates may be optimistic \\
Additional evaluation, not pre-specified & C-MAPSS FD001--FD004 (run once with unchanged code) & Code fixed before the run, but \textit{H} = 20 was designated after these results were known \\
Pre-specified & N-CMAPSS (protocol 2026-09-30), battery (2026-09-30), bearings (2026-10-01): data, model, grids, competitors, splits, levels and four hypotheses & Protocol hashed before the analysis; N-CMAPSS had been used earlier for fixed-time prediction \\
Corrected re-analysis (post hoc) & Battery and bearing analyses with training-only $\kappa$ (Section 5.2), including the $\kappa$ sensitivity & Deviation from the hashed protocols, made after their results were known \\
Exploratory (post hoc) & Fixed-time decomposition; timing controls; health-indicator alarm; battery capacity baselines; identical-protocol competitors and networks with matched inputs; ablation; oracle-timed control; shift study; all cost analyses; other windows; 14-channel N-CMAPSS; STW-T on seen data; thirty random 50-engine test sets; bearing onset diagnosis & Describe the method; not confirmations \\
\bottomrule
\end{tabular}
\end{table*}

The model and grids were developed on PHM08 after fixed-time results of several models were known, and the C-MAPSS datasets were then run once with unchanged code. \textit{H} = 20 was designated afterwards by the criterion of Section 3.1 computed on PHM08 (the early limit could bind for 95.4\% of alarms at \textit{H} = 20, against 65.6\% at \textit{H} = 15). Three protocols were then written in sequence, each after the previous result was known; each fixed the data, $\kappa$ (for aggregated data), the unchanged recipe, two competitors under the same tuning (an age-only regressor, and a multilayer perceptron with a Kalman filter, Section 5.4), five-fold cross-validation by unit (seed 2026), and four hypotheses at \textit{H} = 20: in-window rate $\ge$ 85\%; late-or-missed rate $\le$ 10\%; higher in-window rate than age-only and than the MLP with Kalman filter (exact McNemar test, two-sided \textit{p} < 0.05). The bearing protocol added a fifth: that a theory-guided shift rule (STW-T, which sets the median out-of-fold error to the window centre) would be no more than 5 points worse than the standard rule. A hash shows that a protocol has not changed since it was hashed, not when it was hashed; the dates are attested only by the authors and by file time stamps. N-CMAPSS had been used earlier in the project for fixed-time prediction with other models, before its protocol was written, although no warning rule had been fitted to it. The training-only $\kappa$ of Section 5.2 is a post hoc deviation from the battery and bearing protocols; we report the protocol results unchanged next to the corrected ones and base our interpretation on the corrected ones.

\subsection{Comparisons}

\textbf{Fixed-time references (PHM08).} Eleven fixed-time models (Supplementary Section S1) were trained on 75\% of the training engines of each outer fold and scored at 30, 50, 70 and 90\% of each test engine's life. To separate horizon from calibration, five of them (age only, hidden Markov model, gradient-boosted trees, LSTM and CNN) are also scored with a shift chosen on the remaining 25\% (calibration engines) at the same fraction of life, by the in-window count. The self-timed recipe's own calibrated estimate is also evaluated at the cycle where exactly \textit{H} true cycles remain (oracle-timed control). This control uses the true life and cannot be deployed; it fixes the horizon and isolates the effect of letting the policy choose its moment.

\textbf{Identical-protocol competitors (PHM08).} In the spirit of the three winning PHM08 approaches we re-implemented a similarity method \cite{wang2008}, an LSTM \cite{heimes2008,zheng2017} and a multilayer perceptron followed by a scalar Kalman filter \cite{peel2008}, and added a convolutional network \cite{li2018}. All were run through the same trigger, each with its warning settings chosen on calibration engines (25\% of each outer training part; seed 7) and the shift grid \textminus{}15 to +10. To separate architecture from inputs, an LSTM, a CNN and a hybrid Transformer were also given exactly our inputs and cap set. The re-implementations are not the original code, and only the networks' caps and warning settings were tuned.

\textbf{Timing controls.} (i) A \textbf{constant-estimate control} keeps our alarm times and replaces the reported estimate by one constant per fold, chosen on training out-of-fold alarms. (ii) A \textbf{health-indicator threshold alarm} has no remaining-life regressor: a linear health indicator, in the spirit of Wang et al. \cite{wang2008}, is fitted to separate the first 30 from the last five cycles of each training unit, and the alarm is raised when its smoothed value stays above a threshold $\theta$ for \textit{k} cycles, reporting a constant \textit{c} $\in$ [\textit{H} \textminus{} 6, \textit{H}]; span, debounce, $\theta$ and \textit{c} are chosen by the same nested procedure. (iii) On the batteries, two \textbf{capacity baselines} use the measured discharge capacity directly. The first is a threshold alarm on smoothed capacity, reporting a constant. The second extrapolates a linear capacity trend over the last 10, 20 or 40 blocks to the end-of-life capacity of 0.88 Ah and feeds the result to the same trigger. All settings are chosen on training cells.

\textbf{Ablation (PHM08).} Each variant changes one component and re-tunes the rest: no shift, no cap, cap 125, no smoothing, no debounce, an untuned rule, seven sensors, no age, age only.

\subsection{Cost policies}

All policies use the same outer folds, the same out-of-fold predictions of the training units, and the same lead times and cost ratios. For each ($\ell$, $\rho$) they are tuned on the training units and applied once to the test units.

\begin{itemize}[leftmargin=1.2em,itemsep=2pt]
\item \textbf{STW, window-tuned}: (\textit{c}, \textit{m}, \textit{k}, \textit{s}) chosen by the window count at each \textit{H} $\in$ \{10, 15, 20, 30, 40\} (up to 1,476 settings per level); then \textit{H} chosen by the lowest training cost rate. This policy therefore uses cost, but only to choose among five levels.
\item \textbf{STW, cost-tuned}: \textit{c}, \textit{m}, \textit{k} and a threshold $\tau$ $\in$ \{1, \ldots{}, 60\} on the smoothed estimate chosen to minimise the training cost rate (2,160 settings).
\item \textbf{STW, cost-tuned and regularised}: the same grid, with either a one-standard-error rule (among settings within one bootstrap standard error of the minimum, the one with the largest mean remaining life at the alarm; 200 bootstrap resamples of the training units) or the minimum of the 80th percentile of the bootstrap distribution of the training cost rate.
\item \textbf{Quantile alarm, cost-tuned}: a gradient-boosted 10\% quantile regressor of remaining life (pinball loss, same hyperparameters, target capped at 90); the alarm is raised when its smoothed value stays below $\tau$ for \textit{k} cycles; \textit{m}, \textit{k} and $\tau$ chosen by cost, plainly or with the one-standard-error rule.
\item \textbf{Health-indicator alarm, cost-tuned}, and on the batteries a \textbf{capacity-threshold alarm, cost-tuned} (plain and one-standard-error).
\item \textbf{Optimal age replacement}, with the age chosen on training lives; \textbf{perfect foresight} and \textbf{run-to-failure} as bounds.
\end{itemize}

On the batteries, $\ell$ and the cycles of use are converted with the fold's $\kappa$. The engine results of the window-tuned, cost-tuned, health-indicator and age policies reproduce the earlier analysis of the submitted version unit by unit.

\subsection{Statistics and what the uncertainty covers}

Proportions carry Clopper--Pearson 95\% intervals \cite{clopper1934}. Paired comparisons of in-window rates use the exact McNemar test \cite{mcnemar1947}, with Holm's correction \cite{holm1979} over the seven competitors at each level of the identical-protocol comparison, over the nine ablation variants, and, for the health-indicator alarm, over the seven datasets of the timing-control table; the constant-estimate comparisons are reported unadjusted, since none was significant even before adjustment. The pre-specified hypothesis tests are reported unadjusted, as fixed in the protocols. A non-significant difference is not read as equivalence, since no equivalence margin was specified. Ratios of cost rates carry paired 95\% intervals from 2,000 bootstrap resamples of test units with the fitted policies held fixed. The 16 lead-time and cost-ratio cells of a dataset share the same units and alarms, so they are not independent replications; we report the range of ratios over cells and, descriptively, the number of cells whose interval excludes 1.

These intervals cover the sampling of test units given the fitted policies. They omit the variability from refitting the regressor, from the assignment of units to folds and from the selection of settings. On PHM08 the thirty repeated random 50-engine test sets of Supplementary Section S9 give an indication of that variability. Four leakage checks were run on PHM08: disjoint outer and inner units, causality of the features, a negative control with permuted lives, and exact reproduction of a refitted fold (Supplementary Section S4).

\section{Results}\label{sec:6}

\subsection{Long-range prediction, near-failure estimation and alarm choice}

\begin{figure*}[!tbp]\centering
\includegraphics[width=\textwidth]{img/Fig2_decomposition.png}
\caption{Fixed-time and self-timed evaluation on PHM08 (218 engines, five-fold cross-validation), window [\textminus{}13, +2]. (a) Fixed-time models scored at fractions of each engine's life, on their raw estimates (dashed) and with a shift chosen on calibration engines at the same fraction (solid); the median true remaining life at each fraction is given in brackets. (b) From long-range prediction to the self-timed alarm: the best fixed-time model at half of life, gradient-boosted trees at half and at 90\% of life, the recipe's calibrated estimate when exactly 20 true cycles remain (oracle-timed; not deployable), and the self-timed first warning at \textit{H} = 20, with 95\% Clopper--Pearson intervals.}
\end{figure*}

Fig. 2 separates the four tasks on PHM08. At half of life (median 105 true cycles left), the best of eleven fixed-time models, a stacked hybrid with a tuned shift, placed 33.5\% of predictions in the window, and gradient-boosted trees 20.6\% raw and 20.2\% with a calibration shift. At this horizon the sensors carry little information: the change of a pooled health signal from its early-life level (median 0.23) is comparable to its cycle-to-cycle noise (0.21), whereas 20 cycles before failure it is 1.08. At 90\% of life (median 21 true cycles left, 5--95\%: 15--29) the same model placed 64.7\% in the window raw and 81.2\% (95\% CI 75.4--86.2\%) with a shift learned on calibration engines. The recipe's own estimate, evaluated where exactly 20 true cycles remain, reached 87.6\%, and the self-timed first warning 91.7\% (200/218). Against the calibrated fixed-time model at 90\% of life, the self-timed warning was correct for 35 engines on which the fixed-time model was not, and the reverse held for 12 (exact McNemar \textit{p} = 0.001).

The contrast between 33.5\% and 91.7\% is therefore mostly one of horizon and calibration. Near failure, a calibrated long-range model already meets the window for most engines, so the self-timed policy improves on it by a smaller margin. That margin combines a horizon-focused target, a debounced trigger and the choice of the moment. The oracle-timed control, which fixes the horizon at 20 cycles, placed fewer estimates in the window than the self-timed policy (87.6\% against 91.7\%; 19 against 10 discordant engines, \textit{p} = 0.136). Choosing its own moment, the policy did not do measurably worse than when told the true moment; this is not evidence that it did better. The fixed-time model was also trained on 75\% of the training engines and the recipe on all of them, which favours the recipe in this comparison. These results concern the models and data evaluated; they do not show that fixed-time prediction in general cannot meet such a window.

\subsection{First warnings on the engine datasets}

\begin{table*}[!tbp]
\footnotesize\setlength{\tabcolsep}{4pt}
\caption{First-warning results of the recipe on the development and additional evaluation datasets (five-fold cross-validation by engine, nested selection of the settings). Correct: $\hat{y}$ \textminus{} 2 $\le$ \textit{r} $\le$ $\hat{y}$ + 13 at the alarm.}
\begin{tabular}{@{}>{\raggedright\arraybackslash}p{\dimexpr0.1274\linewidth-2\tabcolsep\relax}>{\raggedright\arraybackslash}p{\dimexpr0.0554\linewidth-2\tabcolsep\relax}>{\raggedright\arraybackslash}p{\dimexpr0.1053\linewidth-2\tabcolsep\relax}>{\raggedright\arraybackslash}p{\dimexpr0.0776\linewidth-2\tabcolsep\relax}>{\raggedright\arraybackslash}p{\dimexpr0.1274\linewidth-2\tabcolsep\relax}>{\raggedright\arraybackslash}p{\dimexpr0.0997\linewidth-2\tabcolsep\relax}>{\raggedright\arraybackslash}p{\dimexpr0.0886\linewidth-2\tabcolsep\relax}>{\raggedright\arraybackslash}p{\dimexpr0.1746\linewidth-2\tabcolsep\relax}>{\raggedright\arraybackslash}p{\dimexpr0.1440\linewidth-2\tabcolsep\relax}@{}}
\toprule
Dataset & \textit{H} & Correct & \% & 95\% CI & Late or missed & Too early & True cycles left at alarm, median (5--95\%) & Alarms with $\ge$ 13 cycles left (\%) \\
\midrule
PHM08 & 10 & 210/218 & 96.3 & 92.9--98.4 & 4 & 4 & 10 (5--17) & 21.7 \\
 & 15 & 205/218 & 94.0 & 90.0--96.8 & 6 & 7 & 15 (8--22) & 65.6 \\
 & 20 & 200/218 & 91.7 & 87.3--95.0 & 8 & 10 & 19 (13--28) & 95.4 \\
 & 30 & 182/218 & 83.5 & 77.9--88.2 & 13 & 23 & 30 (20--43) & 100.0 \\
 & 40 & 173/218 & 79.4 & 73.4--84.5 & 22 & 23 & 38.5 (28--51) & 100.0 \\
FD001 & 10 & 98/100 & 98.0 & 93.0--99.8 & 2 & 0 & 11 (7--17) & 32.0 \\
 & 15 & 95/100 & 95.0 & 88.7--98.4 & 4 & 1 & 14 (8--21) & 59.0 \\
 & 20 & 88/100 & 88.0 & 80.0--93.6 & 4 & 8 & 19 (13--29) & 96.0 \\
FD002 & 10 & 250/260 & 96.2 & 93.0--98.1 & 5 & 5 & 11 (5--17) & 32.7 \\
 & 15 & 252/260 & 96.9 & 94.0--98.7 & 4 & 4 & 15 (9--22) & 73.1 \\
 & 20 & 246/260 & 94.6 & 91.1--97.0 & 7 & 7 & 19 (12--26) & 93.8 \\
FD003 & 10 & 94/100 & 94.0 & 87.4--97.8 & 5 & 1 & 11 (5--17) & 32.7 \\
 & 15 & 94/100 & 94.0 & 87.4--97.8 & 4 & 2 & 16 (8--24) & 77.8 \\
 & 20 & 96/100 & 96.0 & 90.1--98.9 & 1 & 3 & 21 (16--29) & 98.0 \\
FD004 & 10 & 225/249 & 90.4 & 86.0--93.7 & 7 & 17 & 10 (4--20) & 33.3 \\
 & 15 & 214/249 & 85.9 & 81.0--90.0 & 18 & 17 & 14 (6--26) & 61.8 \\
 & 20 & 206/249 & 82.7 & 77.5--87.2 & 17 & 26 & 20 (12--34) & 94.4 \\
\bottomrule
\end{tabular}\par\vspace{3pt}{\raggedright\scriptsize Pooled over PHM08 and FD001--FD004: 877/927 (94.6\%) at \textit{H} = 10, 860/927 (92.8\%) at \textit{H} = 15, 836/927 (90.2\%) at \textit{H} = 20. The five datasets come from one simulator family, so the pooled rate is not an estimate for unrelated systems.\par}
\end{table*}

Table 3 gives the main results. At \textit{H} = 20, 200 of 218 PHM08 engines (91.7\%, 95\% CI 87.3--95.0\%) received a correct first warning; 8 were late, 10 too early and none were missed. The alarm came a median of 19 cycles before failure (5--95\%: 13--28), and 95.4\% of alarms were raised with at least 13 cycles left, so the early limit was active. On C-MAPSS the rate was 88.0\% on FD001, 94.6\% on FD002, 96.0\% on FD003 and 82.7\% on FD004, which has two fault modes and six regimes; 6 of its 26 too-early alarms came more than 30 cycles early. Pooled over the 927 simulated engines it was 90.2\%. On PHM08, lower warning levels gave higher rates but less notice (96.3\% at \textit{H} = 10), and \textit{H} = 30 and 40 gave 83.5\% and 79.4\%; on FD002 and FD003 the rate did not fall monotonically with \textit{H} (Fig. 3). Over thirty random splits of PHM08 into 168 training and 50 test engines, the rate at \textit{H} = 20 averaged 91.8\% (standard deviation 3.3 points; Supplementary Section S9).

\begin{figure}[!tbp]\centering
\includegraphics[width=\linewidth]{img/Fig4_level_tradeoff.png}
\caption{Share of first warnings in the window as a function of the warning level \textit{H} (five-fold cross-validation). Error bars: 95\% Clopper--Pearson intervals for the pre-specified N-CMAPSS analysis.}
\end{figure}

\subsection{What the window score rewards}

\begin{table*}[!tbp]
\footnotesize\setlength{\tabcolsep}{4pt}
\caption{Timing controls at \textit{H} = 20 (five-fold cross-validation): the recipe, the same alarm times with one constant estimate per fold, and a health-indicator threshold alarm without a remaining-life regressor.}
\begin{tabular}{@{}>{\raggedright\arraybackslash}p{\dimexpr0.1883\linewidth-2\tabcolsep\relax}>{\raggedright\arraybackslash}p{\dimexpr0.1496\linewidth-2\tabcolsep\relax}>{\raggedright\arraybackslash}p{\dimexpr0.1496\linewidth-2\tabcolsep\relax}>{\raggedright\arraybackslash}p{\dimexpr0.1773\linewidth-2\tabcolsep\relax}>{\raggedright\arraybackslash}p{\dimexpr0.1662\linewidth-2\tabcolsep\relax}>{\raggedright\arraybackslash}p{\dimexpr0.1691\linewidth-2\tabcolsep\relax}@{}}
\toprule
Dataset & Recipe: correct (\%) & Constant estimate: correct (\%) & Discordant (recipe only / control only); \textit{p} & Health-indicator alarm: correct (\%) & Discordant; \textit{p} (Holm) \\
\midrule
PHM08 & 200/218 (91.7) & 200 (91.7) & 7 / 7; 1.000 & 122 (56.0) & 83 / 5; < 0.001 \\
FD001 & 88/100 (88.0) & 90 (90.0) & 2 / 4; 0.688 & 51 (51.0) & 41 / 4; < 0.001 \\
FD002 & 246/260 (94.6) & 243 (93.5) & 9 / 6; 0.607 & 142 (54.6) & 105 / 1; < 0.001 \\
FD003 & 96/100 (96.0) & 94 (94.0) & 4 / 2; 0.688 & 57 (57.0) & 40 / 1; < 0.001 \\
FD004 & 206/249 (82.7) & 215 (86.3) & 5 / 14; 0.064 & 123 (49.4) & 94 / 11; < 0.001 \\
N-CMAPSS & 68/99 (68.7) & 70 (70.7) & 7 / 9; 0.804 & 53 (53.5) & 26 / 11; 0.020 \\
Battery & 118/123 (95.9) & 115 (93.5) & 7 / 4; 0.549 & 67 (54.5) & 54 / 3; < 0.001 \\
Engines pooled (PHM08, FD001--FD004) & 836/927 (90.2) & 842 (90.8) & 27 / 33; 0.519 & 495 (53.4) & -- \\
\bottomrule
\end{tabular}\par\vspace{3pt}{\raggedright\scriptsize Exact McNemar tests. Constant-estimate column: unadjusted \textit{p}; health-indicator column: Holm-adjusted over the seven datasets. Battery: blocks of the training-only $\kappa$ of each fold (rerun for this revision). All levels \textit{H} = 10--40 are in Supplementary Section S10.\par}
\end{table*}

Because the alarm is raised when the smoothed, shifted estimate first stays below \textit{H}, the reported estimate varies little between units: on PHM08 at \textit{H} = 20 its standard deviation at the alarm was 2.1 cycles (5--95\%: 11.2--18.2), against 4.8 cycles for the true remaining life. Correctness, $\hat{y}$ \textminus{} 2 $\le$ \textit{r} $\le$ $\hat{y}$ + 13, therefore depends mainly on where the alarm falls. For $\hat{y}$ between 11.2 and 18.2 the accepted range of \textit{r} runs from at least 9.2 to at most 31.2 cycles; a fixed band is only an approximation of this condition. Table 4 tests this reading. With our alarm times and one constant estimate per fold, the in-window rate was 90.8\% over the five engine datasets, against 90.2\% with the model's estimate (27 against 33 discordant engines, \textit{p} = 0.519). The constant was numerically lower on FD002, FD003, Battery, higher on FD001, FD004, N-CMAPSS, and equal on PHM08. Per dataset the constant scored between \textminus{}2.4 and +3.6 percentage points relative to the model's estimate. No difference was significant, but no equivalence margin was tested either, so the right conclusion is limited: the model's estimate added no detectable information beyond the alarm time.

The regressor is not therefore dispensable. The health-indicator threshold alarm, which times its alarm on a classical linear health indicator, placed 49.4--57.0\% of first warnings in the window, significantly fewer than the recipe on every dataset after Holm's correction (largest adjusted \textit{p} = 0.020, N-CMAPSS). Many of its alarms were too early (64 of 218 on PHM08). In this recipe, what the regressor contributes is a health indicator calibrated in cycles to failure, so that one threshold fires at a similar remaining life on units that degrade at different rates. The window score rewards that timing, not the precision of the number attached to the alarm.

\begin{figure*}[!tbp]\centering
\includegraphics[width=\textwidth]{img/Fig5_error_leadtime.png}
\caption{Distribution of (a) the warning error and (b) the true remaining life at the alarm at \textit{H} = 20 for PHM08 (five-fold cross-validation) and N-CMAPSS (pre-specified analysis). Shaded: the window; dotted: 13 cycles, below which the early limit practically cannot bind.}
\end{figure*}

\subsection{Matched comparisons and ablation}

\begin{table*}[!tbp]
\footnotesize\setlength{\tabcolsep}{4pt}
\caption{Identical first-warning protocol on PHM08 (218 engines; fit on 75\% of each outer training part, warning settings chosen on the other 25\%). Upper rows of each level: re-implemented approaches with their original inputs; last three rows: networks given exactly the inputs and cap set of the recipe.}
\begin{tabular}{@{}>{\raggedright\arraybackslash}p{\dimexpr0.2493\linewidth-2\tabcolsep\relax}>{\raggedright\arraybackslash}p{\dimexpr0.0499\linewidth-2\tabcolsep\relax}>{\raggedright\arraybackslash}p{\dimexpr0.1274\linewidth-2\tabcolsep\relax}>{\raggedright\arraybackslash}p{\dimexpr0.1163\linewidth-2\tabcolsep\relax}>{\raggedright\arraybackslash}p{\dimexpr0.0942\linewidth-2\tabcolsep\relax}>{\raggedright\arraybackslash}p{\dimexpr0.0886\linewidth-2\tabcolsep\relax}>{\raggedright\arraybackslash}p{\dimexpr0.1551\linewidth-2\tabcolsep\relax}>{\raggedright\arraybackslash}p{\dimexpr0.1192\linewidth-2\tabcolsep\relax}@{}}
\toprule
Method & \textit{H} & Correct (\%) & 95\% CI & Late or missed & Too early & Discordant (ours only / competitor only) & Holm-adjusted \textit{p} \\
\midrule
Recipe (ours) & 10 & 206 (94.5) & 90.6--97.1 & 9 & 3 & -- & -- \\
MLP + Kalman filter (Peel-style) & 10 & 195 (89.4) & 84.6--93.2 & 9 & 14 & 22 / 11 & 0.240 \\
LSTM (Heimes-style) & 10 & 173 (79.4) & 73.4--84.5 & 16 & 29 & 42 / 9 & < 0.001 \\
Similarity (Wang-style) & 10 & 160 (73.4) & 67.0--79.1 & 21 & 37 & 55 / 9 & < 0.001 \\
CNN \cite{li2018} & 10 & 138 (63.3) & 56.5--69.7 & 40 & 40 & 78 / 10 & < 0.001 \\
LSTM, our inputs and cap set & 10 & 201 (92.2) & 87.8--95.4 & 10 & 7 & 11 / 6 & 0.665 \\
CNN, our inputs and cap set & 10 & 179 (82.1) & 76.4--87.0 & 20 & 19 & 35 / 8 & < 0.001 \\
Hybrid Transformer, our inputs and cap set & 10 & 206 (94.5) & 90.6--97.1 & 10 & 2 & 8 / 8 & 1.000 \\
Recipe (ours) & 20 & 195 (89.4) & 84.6--93.2 & 10 & 13 & -- & -- \\
MLP + Kalman filter (Peel-style) & 20 & 180 (82.6) & 76.9--87.4 & 15 & 23 & 30 / 15 & 0.071 \\
LSTM (Heimes-style) & 20 & 140 (64.2) & 57.5--70.6 & 29 & 49 & 70 / 15 & < 0.001 \\
Similarity (Wang-style) & 20 & 138 (63.3) & 56.5--69.7 & 34 & 46 & 64 / 7 & < 0.001 \\
CNN \cite{li2018} & 20 & 121 (55.5) & 48.6--62.2 & 46 & 51 & 82 / 8 & < 0.001 \\
LSTM, our inputs and cap set & 20 & 176 (80.7) & 74.9--85.7 & 27 & 15 & 36 / 17 & 0.038 \\
CNN, our inputs and cap set & 20 & 172 (78.9) & 72.9--84.1 & 22 & 24 & 40 / 17 & 0.013 \\
Hybrid Transformer, our inputs and cap set & 20 & 201 (92.2) & 87.8--95.4 & 5 & 12 & 13 / 19 & 0.377 \\
\bottomrule
\end{tabular}\par\vspace{3pt}{\raggedright\scriptsize Exact McNemar tests against the recipe; Holm correction over the seven competitors at each level. The recipe scores slightly lower here than in Table 3 because it is fitted on 75\% of the training engines.\par}
\end{table*}

Under the identical protocol (Table 5), the recipe placed 89.4\% of PHM08 first warnings in the window at \textit{H} = 20. The LSTM, similarity and CNN re-implementations with their original inputs placed 64.2\%, 63.3\% and 55.5\%, significantly fewer at both levels (Holm-adjusted \textit{p} < 0.001). The MLP with Kalman filter, which shares our feature encoder, placed 82.6\%. The difference was not significant after correction (30 against 15 discordant engines, Holm-adjusted \textit{p} = 0.071 at \textit{H} = 20; 0.240 at \textit{H} = 10). Given our inputs and cap set, the LSTM and CNN rose to 80.7\% and 78.9\%, still significantly below the recipe at \textit{H} = 20. The hybrid Transformer reached 92.2\%, numerically higher than the recipe (13 against 19 discordant engines, \textit{p} = 0.377), at about 36 minutes of training per fold against about 34 seconds. Much of the gap to the original re-implementations is therefore due to inputs and targets. Since only caps and warning settings of the networks were tuned, these comparisons do not rank architectures.

The ablation (Fig. 5b) reproduced the main result with the full recipe (200/218). Among the trigger and target components, removing the shift caused the largest loss (77.1\%, with 49 late alarms; Holm-adjusted \textit{p} < 0.001). Training without a cap gave 81.7\% (\textit{p} = 0.002), and an untuned rule gave 78.4\%. Seven sensors gave 83.9\%, sensor features without age 85.8\%, and age alone 14.2\%. The conventional cap of 125 (89.9\%), no smoothing (90.4\%) and no debounce (89.4\%) were not significantly different from the full recipe. With 218 engines, such tests can miss differences of a few percentage points.

\begin{figure*}[!tbp]\centering
\includegraphics[width=\textwidth]{img/Fig6_comparison_ablation.png}
\caption{(a) In-window rate under the identical protocol of Table 5 with 95\% intervals. (b) Component ablation on PHM08 at \textit{H} = 20 (nested protocol); each bar changes one component and re-tunes the remaining settings.}
\end{figure*}

\subsection{The shift moves the alarm}

Proposition 1 assumes that the error distribution does not change with the shift. To see how far this holds, we varied the shift on the test units of each fold by up to $\pm$4 around the value selected on training data at \textit{H} = 20, keeping the other settings, and recorded the alarm cycle, the true remaining life at the alarm and the unshifted error $\varepsilon$ (Fig. 6). On every dataset the shift acted mainly by moving the alarm. Each unit of shift changed the median true remaining life at the alarm by \textminus{}1.0 to \textminus{}0.9 cycles (blocks for the battery), against \textminus{}0.1 to 0.1 for the median of $\varepsilon$. Assumption A1 is therefore violated by construction, as the stopping time changes, yet the error at the alarm changed slowly.

Within $\pm$2 of the selected shift, the largest difference between the realised in-window rate and the rate predicted by Proposition 1 from the errors at the selected shift was 1.6 to 8.0 percentage points, depending on the dataset (largest on FD001). Within $\pm$4 it was 4.2 to 8.0 points. On the test units, the best shift placed up to 5.0 points more first warnings in the window than the shift selected on training data (FD001), a comparison that favours the test optimum. On these data the proposition therefore describes the direction and rough size of the response to the shift, with errors of several points. It is not a bound for the sequential policy, and it says nothing about what other features or models could achieve. The submitted version's statements that no choice of shift could rescue a low rate, and that the concentration of the error could serve as a check before deployment, are withdrawn.

\begin{figure*}[!tbp]\centering
\includegraphics[width=\textwidth]{img/Fig8_shift_A1.png}
\caption{Assumption A1 of Proposition 1, test units at \textit{H} = 20 with the other settings fixed. (a) Realised in-window rate (solid) and the rate predicted by Proposition 1 from the errors at the selected shift (dashed) as the shift moves away from the value selected on training data. (b) Change of the median true remaining life at the alarm relative to the selected shift; dotted: one cycle less per unit of shift. Battery in blocks of the training-only $\kappa$.}
\end{figure*}

\subsection{Maintenance cost}

\begin{table*}[!tbp]
\footnotesize\setlength{\tabcolsep}{4pt}
\caption{Cost of the window-tuned alarms relative to each alternative policy, by dataset. Each entry: median (minimum--maximum) over the 16 lead-time and cost-ratio cells of the ratio of cost rates (window-tuned / alternative; below 1 favours the window-tuned policy); then the number of cells in which the paired 95\% bootstrap interval lies entirely below 1 / entirely above 1.}
\begin{tabular}{@{}>{\raggedright\arraybackslash}p{\dimexpr0.1199\linewidth-2\tabcolsep\relax}>{\raggedright\arraybackslash}p{\dimexpr0.1253\linewidth-2\tabcolsep\relax}>{\raggedright\arraybackslash}p{\dimexpr0.1253\linewidth-2\tabcolsep\relax}>{\raggedright\arraybackslash}p{\dimexpr0.1253\linewidth-2\tabcolsep\relax}>{\raggedright\arraybackslash}p{\dimexpr0.1253\linewidth-2\tabcolsep\relax}>{\raggedright\arraybackslash}p{\dimexpr0.1253\linewidth-2\tabcolsep\relax}>{\raggedright\arraybackslash}p{\dimexpr0.1253\linewidth-2\tabcolsep\relax}>{\raggedright\arraybackslash}p{\dimexpr0.1282\linewidth-2\tabcolsep\relax}@{}}
\toprule
Dataset & Age replacement & Health-indicator alarm & STW, cost-tuned & STW, cost-tuned, 1-SE & STW, cost-tuned, 80th pct. & Quantile alarm, cost-tuned & Quantile alarm, 1-SE \\
\midrule
PHM08 & 0.60 (0.45--0.64); 16/0 & 0.96 (0.79--0.99); 10/0 & 0.93 (0.61--0.99); 0/0 & 1.01 (0.84--1.22); 1/4 & 0.93 (0.61--0.99); 0/0 & 1.00 (0.63--1.05); 0/3 & 1.00 (0.82--1.21); 0/2 \\
FD001 & 0.59 (0.45--0.65); 16/0 & 0.83 (0.49--0.95); 9/0 & 0.96 (0.68--1.03); 0/4 & 1.00 (0.99--1.02); 4/7 & 0.96 (0.68--1.03); 0/4 & 1.04 (1.03--1.06); 0/16 & 1.02 (1.01--1.04); 0/16 \\
FD002 & 0.65 (0.55--0.69); 15/0 & 0.94 (0.80--1.16); 12/0 & 0.97 (0.66--1.06); 0/0 & 1.01 (0.85--1.05); 1/11 & 1.00 (0.66--1.06); 0/4 & 1.01 (0.74--1.06); 0/4 & 1.02 (0.99--1.24); 1/14 \\
FD003 & 0.55 (0.42--0.61); 16/0 & 0.85 (0.63--0.95); 16/0 & 1.01 (0.51--1.03); 0/8 & 0.98 (0.98--0.99); 16/0 & 1.01 (0.51--1.03); 0/8 & 0.90 (0.51--1.02); 0/4 & 0.98 (0.97--1.00); 12/3 \\
FD004 & 0.56 (0.48--0.67); 15/0 & 0.97 (0.84--1.17); 2/0 & 0.97 (0.71--1.01); 0/0 & 1.00 (0.86--1.19); 0/3 & 0.97 (0.71--1.01); 0/0 & 1.02 (0.76--1.20); 0/3 & 1.04 (1.01--1.41); 0/11 \\
N-CMAPSS & 0.98 (0.69--1.06); 0/0 & 0.80 (0.44--1.05); 6/0 & 1.17 (0.79--1.73); 0/9 & 1.14 (0.78--1.73); 0/9 & 1.16 (0.79--1.73); 0/9 & 1.12 (0.78--1.76); 0/6 & 1.11 (0.77--1.73); 0/4 \\
Battery & 0.49 (0.33--0.52); 16/0 & 0.98 (0.91--1.06); 11/4 & 1.03 (0.98--1.06); 1/13 & 0.98 (0.94--1.02); 12/2 & 1.02 (0.98--1.06); 1/12 & 1.00 (0.97--1.05); 2/6 & 0.96 (0.91--1.00); 15/0 \\
\bottomrule
\end{tabular}\par\vspace{3pt}{\raggedright\scriptsize Five-fold cross-validation; lead times $\ell$ $\in$ \{0, 5, 10, 15\} cycles (battery 0, 15, 30, 45 cycles) and failure-cost ratios $\rho$ $\in$ \{5, 10, 20, 50\}. The 16 cells of a dataset share units and alarms and are not independent; the counts are descriptive. Battery: training-only $\kappa$. The battery capacity-threshold policies and the cost rates relative to perfect foresight for every cell are in Supplementary Section S11.\par}
\end{table*}

Table 6 and Fig. 7 compare the policies under identical splits, information, lead times and cost ratios. Against optimal age replacement, the window-tuned alarms were cheaper with a paired interval excluding equality in 16, 16, 15, 16, 15, 16 of 16 cells on PHM08, FD001, FD002, FD003, FD004, Battery, respectively. Their cost rate was 1.07--1.28 times that of perfect foresight on PHM08, against 1.70--2.37 for age replacement. N-CMAPSS was the exception: in no cell did the interval exclude equality, and the point estimate favoured age replacement in 4 of 16 cells.

Against policies tuned directly to cost the results depend on the dataset (Table 6). Plain cost tuning of the same regressor was significantly cheaper than window tuning in 4, 8, 9 and 13 of the 16 cells on FD001, FD003, N-CMAPSS and the batteries (none on PHM08, FD002 and FD004), and significantly more expensive in 1 of the 16 cells on the batteries (none elsewhere). Its occasional large losses, which the submitted version emphasised, remain. On PHM08 with $\ell$ = 10 and $\rho$ = 50 it cost 1.75 times perfect foresight, against 1.10 for window tuning, with 3 test failures against 0. The one-standard-error rule reduced that loss to 1.30, still above window tuning there. Across cells it was significantly cheaper than window tuning in 4, 7, 11, 3, 9 and 2 of the 16 cells on PHM08, FD001, FD002, FD004, N-CMAPSS and the batteries (none on FD003), and significantly more expensive in 1, 4, 1, 16 and 12 of the 16 cells on PHM08, FD001, FD002, FD003 and the batteries (none on FD004 and N-CMAPSS). The cost-tuned 10\% quantile alarm was significantly cheaper than window tuning in 3, 16, 4, 4, 3, 6 and 6 of the 16 cells on PHM08, FD001, FD002, FD003, FD004, N-CMAPSS and the batteries, and more expensive in 2 of the 16 cells on the batteries (none elsewhere). The health-indicator alarm was usually the more expensive: window tuning was significantly cheaper in 10, 9, 12, 16, 2, 6 and 11 of the 16 cells on PHM08, FD001, FD002, FD003, FD004, N-CMAPSS and the batteries, and more expensive in 4 of the 16 cells on the batteries (none elsewhere). On the batteries, the cost-tuned capacity-threshold policy was significantly cheaper than window tuning in 4 cells and never significantly more expensive.

Relative to the cheapest condition-based policy in each cell, the mean excess cost over the 112 cells was 7.9\% for window tuning, 13.6\% for plain cost tuning, 5.6\% with the one-standard-error rule, 10.2\% for the quantile alarm (5.9\% with the one-standard-error rule), 26.3\% for the health-indicator alarm and 85.6\% for age replacement. These means are descriptive, because the cells are correlated. The policy with the lowest mean excess cost differed between datasets: window-tuned STW on PHM08, cost-tuned quantile alarm on FD001, quantile alarm with the 1-SE rule on FD002, window-tuned STW on FD003, quantile alarm with the 1-SE rule on FD004, cost-tuned STW (80th percentile) on N-CMAPSS, capacity threshold with the 1-SE rule on Battery. The largest single excess cost was 94.8\% for plain cost tuning (FD003), 75.7\% for window tuning (N-CMAPSS, $\ell$ = 5, $\rho$ = 50), 53.0\% with the one-standard-error rule (Battery) and 102.1\% for the quantile alarm (N-CMAPSS). Window tuning thus avoided the largest losses of plain cost tuning on several datasets, plausibly through the margin that the narrow late side of the window builds in, but that mechanism was not tested and it did not hold on N-CMAPSS. No tuning rule was cheapest on every dataset, and window tuning is not a reliable surrogate for cost optimisation.

\begin{figure*}[!tbp]\centering
\includegraphics[width=\textwidth]{img/FigR3_cost_audit.png}
\caption{Long-run cost rate relative to perfect foresight (1.00 = replacement one cycle before every failure) for policies tuned on training units: (a) without lead time and $\rho$ = 10; (b) with a lead time of 10 cycles (30 cycles for the battery) and $\rho$ = 50. Five-fold cross-validation; battery with training-only $\kappa$. Off-scale values are labelled.}
\end{figure*}

\subsection{Retraining on new data}

\begin{table*}[!tbp]
\footnotesize\setlength{\tabcolsep}{4pt}
\caption{Retraining the recipe on new data (\textit{H} = 20 blocks or cycles; five-fold cross-validation by unit unless stated). Protocol: block length $\kappa$ as in the hashed protocol, from all units; corrected: $\kappa$ from the training units of each split (Section 5.2).}
\begin{tabular}{@{}>{\raggedright\arraybackslash}p{\dimexpr0.2414\linewidth-2\tabcolsep\relax}>{\raggedright\arraybackslash}p{\dimexpr0.1312\linewidth-2\tabcolsep\relax}>{\raggedright\arraybackslash}p{\dimexpr0.1050\linewidth-2\tabcolsep\relax}>{\raggedright\arraybackslash}p{\dimexpr0.1050\linewidth-2\tabcolsep\relax}>{\raggedright\arraybackslash}p{\dimexpr0.0997\linewidth-2\tabcolsep\relax}>{\raggedright\arraybackslash}p{\dimexpr0.0892\linewidth-2\tabcolsep\relax}>{\raggedright\arraybackslash}p{\dimexpr0.0945\linewidth-2\tabcolsep\relax}>{\raggedright\arraybackslash}p{\dimexpr0.1339\linewidth-2\tabcolsep\relax}@{}}
\toprule
Dataset and analysis & Recipe: correct (\%) & 95\% CI & Late / early / missed & MLP + Kalman filter (\%) & Age-only (\%) & Capacity trend (\%) & True life left at alarm, median \\
\midrule
N-CMAPSS, 99 engines (pre-specified) & 68/99 (68.7) & 58.6--77.6 & 18 / 8 / 5 & 58.6 & 24.2 & -- & 17 cycles \\
Battery, 123 cells, protocol ($\kappa$ = 3) & 120/123 (97.6) & 93.0--99.5 & 2 / 0 / 1 & 85.4 & 0.8 & -- & 57 cycles \\
Battery, 123 cells, corrected ($\kappa$ $\in$ \{3, 4\}) & 118/123 (95.9) & 90.8--98.7 & 2 / 2 / 1 & 80.5 & 2.4 & 100.0 & 62 cycles \\
Battery, batch held out, protocol ($\kappa$ = 3) & 86/123 (69.9) & 61.0--77.9 & 27 / 10 / 0 & 39.0 & 3.3 & -- & 57 cycles \\
Battery, batch held out, corrected & 80/123 (65.0) & 55.9--73.4 & 30 / 9 / 4 & 50.4 & 3.3 & 99.2 & 52 cycles \\
Battery, 119 cells, corrected & 117/119 (98.3) & 94.1--99.8 & 2 / 0 / 0 & 82.4 & 3.4 & -- & 60 cycles \\
Bearings, 17, protocol ($\kappa$ = 7) & 0/17 (0.0) & 0.0--19.5 & 0 / 7 / 10 & 0.0 & 5.9 & -- & -- \\
Bearings, 17, corrected & 0/17 (0.0) & 0.0--19.5 & 0 / 7 / 10 & 0.0 & 5.9 & -- & -- \\
Bearings, official split (11 test), corrected & 1/11 (9.1) & 0.2--41.3 & 0 / 2 / 8 & 9.1 & 0.0 & -- & -- \\
\bottomrule
\end{tabular}\par\vspace{3pt}{\raggedright\scriptsize Protocol hypotheses at \textit{H} = 20 (in-window $\ge$ 85\%; late or missed $\le$ 10\%; better than age-only; better than the MLP with Kalman filter): N-CMAPSS, only the third met; battery cross-validation, all four met under both $\kappa$ rules; bearings, none met under either rule. Capacity trend: linear extrapolation of measured capacity to the end-of-life level of 0.88 Ah with the same trigger (exploratory). Battery leave-one-batch-out: inner seed 100 + held-out batch.\par}
\end{table*}

Table 7 and Fig. 8 summarise the three datasets on which the recipe was retrained under protocols fixed in advance. No model was transferred between them.

\textbf{Battery cells.} With $\kappa$ from the training cells, the recipe placed 118 of 123 first warnings in the window under cross-validation by cell (95.9\%, 95\% CI 90.8--98.7\%; protocol analysis 97.6\%), with 2 late, 2 too early and 1 missed. All four protocol criteria were met: the recipe was correct on 20 cells on which the MLP with Kalman filter was not, and the reverse held for 1 (\textit{p} < 0.001). Age-only placed 2.4\%, because cell lives range from 148 to 1935 cycles. When a whole batch was held out, the rate fell to 65.0\% (80/123; protocol analysis 69.9\%), with 30 late and 4 missed cells. It was still above the MLP with Kalman filter (50.4\%; \textit{p} = 0.030), but not on batch 2. That batch has the shortest lives, and a model trained on the longer-lived batches 1 and 3 ($\kappa$ = 4) expected more life than its cells had: 15 of 43 correct, 26 late, against 27 for the MLP with Kalman filter. Batch 2 is also where the correction of $\kappa$ mattered: with $\kappa$ = 3, as in the protocol, 24 of its 43 cells were correct. The held-out-batch rate therefore depends on the time scale as well as on the shift in conditions, and the $\kappa$ sensitivity analysis below covers cross-validation only. The within-batch result shows what the recipe achieves under familiar test conditions. The held-out-batch result is the relevant estimate for a new campaign with different protocols and rest periods, and it falls well short of the 85\% criterion.

A battery-specific baseline puts both numbers in context. Battery end of life is \textit{defined} by discharge capacity falling to 0.88 Ah, and capacity is one of the measured inputs. Linearly extrapolating the recent capacity trend to that level, with the same trigger, placed 123/123 (100.0\%) of first warnings in the window within batches and 122/123 (99.2\%) with a batch held out, significantly more than the recipe in the second case (43 against 1 discordant cells, \textit{p} < 0.001). Its alarms came later: with a batch held out, a median of 44 cycles before failure against 52 for the recipe; within batches, only 60.2\% of its alarms were raised early enough to be too early, against 95.1\% for the recipe. A threshold on smoothed capacity reached only 64.2\% and 49.6\%. On this score, then, a simple physical extrapolation did at least as well as the learned recipe on these cells (within batches 5 against 0 discordant cells, \textit{p} = 0.063). Much of the battery success reflects a failure criterion that is a threshold on a measured quantity, not an ability peculiar to the recipe.

With the physical task fixed (\textit{H} = 60 cycles, window [\textminus{}39, +6] cycles), block lengths $\kappa$ $\in$ \{1, 2, 3, 4, 5, 6\} gave 98.4, 95.1, 98.4, 95.1, 97.6, 98.4\% under cross-validation (Supplementary Table S3), so the within-batch result does not hinge on the aggregation.

\textbf{N-CMAPSS.} The recipe placed 68 of 99 first warnings in the window (68.7\%, 95\% CI 58.6--77.6\%); 18 were late and 5 engines were missed. It was clearly better than age-only (24.2\%; criterion 3 met), but neither threshold was met, and its advantage over the MLP with Kalman filter (58.6\%) was not significant (25 against 15 discordant engines, \textit{p} = 0.154). Lives are short, so a 20-cycle alarm comes at roughly three quarters of life, and the true remaining life at the alarm was spread more widely than on PHM08 (Fig. 4b). An exploratory re-run with all 14 channels did not help (Supplementary Section S6).

\textbf{Bearings.} The recipe failed. None of the 17 first warnings fell in the window, in the protocol analysis or with $\kappa$ from the training bearings (7 too early and 10 missed in the corrected analysis). The competitors were also almost always wrong (age-only 1, MLP with Kalman filter 0 of 17). With the physical task fixed (\textit{H} = 140 snapshots), block lengths of 3, 5, 7, 10, 14 snapshots gave at most 1 of 17. The pre-specified fifth hypothesis, that STW-T would be no more than 5 points worse, was met only trivially (0\% against 0\%), which cannot establish non-inferiority. A post hoc diagnosis with an ad hoc criterion (Supplementary Section S7) found that in 11 of 17 bearings the vibration level rose to and stayed above 1.5 times its early-life level only within the last six blocks before failure (about seven minutes), or never. 6 of these were never warned, which is consistent with failure becoming visible too late for a 20-block alarm. The other 5 received alarms 63--211 blocks before failure, long before any sustained rise, which late visibility cannot explain. Of the 6 bearings with longer degradation, 4 were missed and 2 warned too early. Late onset is therefore at most a partial explanation. Small training sets (13--14 bearings per fold), three operating conditions and the feature set are competing explanations that were not tested. The result shows that this recipe, with these features, did not work on these bearings; it does not show that advance warning of bearing failure is impossible.

\begin{figure*}[!tbp]\centering
\includegraphics[width=\textwidth]{img/Fig7_battery_v3.png}
\caption{Battery cells, block length $\kappa$ from the training cells of each split. (a) Discharge capacity of the 123 cells (grey) and the first warning of the recipe at \textit{H} = 20 blocks under cross-validation by cell, coloured by outcome; the missed cell has no marker. (b) In-window rate under cross-validation by cell and with a whole batch held out, for the recipe, the two pre-specified competitors and two capacity baselines.}
\end{figure*}

\section{Discussion}\label{sec:7}

\subsection{When did first-warning policies give useful decisions?}

The evidence answers the central question only conditionally, and the conditions below are a post hoc reading of eight datasets, not a tested rule. The recipe produced alarms that met the window and were cheaper than age replacement on the datasets where two things held: the approach of failure became visible, consistently across units, some time before the notice required; and the units to be monitored resembled those used for training. These were the five simulated engine datasets and the battery cells within batches. Under a change of test campaign (a held-out battery batch) the rate fell from 95.9\% to 65.0\%. Where the remaining life at the alarm was widely spread (N-CMAPSS), the rate was 68.7\%, and its alarms were not significantly cheaper than age replacement in any cost setting. On the bearings, with the features used, the recipe failed, and late visibility of failure explains at most part of that failure.

Three qualifications matter for practice. (i) Under the asymmetric window, the score measures \textit{when} the alarm is raised. The reported estimate was nearly constant at the alarm, and a constant in its place scored within a few percentage points on every dataset (no equivalence margin was tested). An operator who adopts such a window as an acceptance criterion is therefore constraining the remaining life at the alarm, roughly to a band of 15 cycles placed by the warning level, and not the precision of the number shown. (ii) A good first-warning rate does not mean that the model is a good long-range predictor, and a poor long-range rate does not mean that a model is poor near failure (Section 6.1). The four tasks need separate evidence, and planning models that set maintenance slots well ahead remain a separate problem \cite{kamariotis2024,mitici2023}. (iii) Domain baselines can be hard to beat. On the battery cells, an extrapolation of measured capacity to the capacity that defines end of life did at least as well as the learned recipe on the window score, and better with a batch held out.

\subsection{Window tuning and cost tuning}

Kamariotis et al. \cite{kamariotis2024} argue that prognostic algorithms and the policies built on them should be judged and tuned by the long-run cost of the decisions they induce. Our results support evaluating alarms by cost, with a qualification about how the cost is minimised. Minimising the empirical cost rate over about two thousand settings sometimes selected alarms with too little margin and produced large losses when failures were expensive. A one-standard-error rule reduced the largest losses, a lower quantile was often cheaper but did not avoid them, and no tuning rule was cheapest on every dataset. Window tuning was cheaper than the one-standard-error rule on FD003 and the batteries and more expensive on several other datasets. Its settings other than the warning level do not depend on cost ratios, which are often uncertain in practice. The window is best regarded as an operational specification, not a proxy for cost. Where cost ratios and lead times can be estimated, tuning to the cost with a safeguard against selection noise should be compared with it on the fleet's own history.

\subsection{Limitations}

\begin{itemize}[leftmargin=1.2em,itemsep=2pt]
\item \textbf{No field data.} Six of the eight datasets are simulations from one simulator family (C-MAPSS), and the battery and bearing data are laboratory measurements; on the batteries, end of life is a threshold on a measured quantity. First-warning rates at \textit{H} = 20 ranged from 0\% to 96\% across datasets and fell under a change of test campaign, so they must be re-estimated for each application.
\item \textbf{Development and selection.} The recipe and grids were developed on PHM08, and the primary level was designated after the C-MAPSS results were known. Only three datasets had protocols fixed in advance; the hashes were not registered, so their timing rests on the authors' records. The training-only $\kappa$ is a post hoc correction, and all controls, baselines, cost analyses and the shift study were added after the main results were known.
\item \textbf{Competitors.} The challenge approaches are re-implementations, not the original code. Only the networks' caps and warning settings were tuned, and the health-indicator and capacity baselines are single representatives of their classes.
\item \textbf{Statistics.} Intervals are conditional on the fitted policies; refitting and fold-assignment variability is reflected only in the repeated PHM08 splits. Non-significant differences are not evidence of equivalence. The many comparisons are adjusted within tables, not across the paper, and cost cells within a dataset are correlated.
\item \textbf{Cost model.} The renewal--reward model assumes a fixed lead time, independent units and replacement as good as new. The cost ratios are illustrative, and the model ignores spare parts, downtime and the grouping of maintenance across units \cite{depater2022,mitici2023}. The window-tuned and cost-tuned policies select on the same out-of-fold predictions that tune their other settings.
\item \textbf{Theory.} Proposition 1 holds for a fixed error distribution. For the sequential policy it is only an approximate description on these data, with errors of up to 8.0 points, not a guarantee.
\end{itemize}

\section{Future work}\label{sec:8}

Four directions follow directly from the limitations. \textbf{Field data and changing conditions}: the battery results under a held-out batch show that alarms tuned on one campaign lose accuracy on another. Recalibrating the shift on the first failures of a new population, or attaching distribution-free uncertainty to the alarm time, should be tested on fleet data with censoring. \textbf{Maintenance policies}: the warning level should be chosen jointly with the order lead time, and the alarms of several units grouped into shared visits \cite{depater2022,mitici2023}. The comparison of Section 6.6 suggests combining the margin that window tuning builds in with explicit cost criteria and probabilistic prognostics, rather than choosing one. \textbf{Test and inspection policies}: for units that are proof-tested or inspected at intervals \cite{dejonge2020}, an alarm could trigger an additional test rather than a replacement, which requires an encoder for sparse observations and a cost model with test costs. \textbf{Fault detection}: for failures that become visible only shortly before they happen, the first-warning rule should be compared with sequential detectors such as cumulative-sum and exponentially weighted moving-average charts \cite{page1954,roberts1959} under the same leakage-safe protocol, and a two-stage design, in which a detector declares the onset of a defect and a remaining-life model then times the replacement, remains an untested hypothesis for the bearings.

\section{Conclusion}\label{sec:9}

Scoring remaining-life models at the moment they raise an alarm changes what is measured. On simulated engines, a simple self-timed recipe placed 91.7\% of PHM08 and 90.2\% of 927 first warnings within an asymmetric window at a 20-cycle warning level, whereas fixed-time models at half of life reached at most 33.5\%. Most of that contrast reflects horizon and calibration: at 90\% of life a calibrated fixed-time model reached 81.2\%. Three distinctions follow. \textit{Prediction accuracy}: the window score says little about it, because a constant estimate at the same alarm times scored within a few percentage points of the model's estimate. \textit{Alarm timing}: this is what the score measures, and at \textit{H} = 20 a regressor calibrated in cycles to failure timed alarms much better than a health-indicator threshold on every dataset. \textit{Maintenance cost}: window-tuned alarms were cheaper than age replacement on every dataset except N-CMAPSS, but alarms tuned directly to cost were cheaper on some datasets and more expensive on others, and no tuning rule was best everywhere, so the window is not a reliable surrogate for cost. Retrained on new data, the recipe reached 95.9\% within battery batches but 65.0\% with a batch held out, where a capacity-trend extrapolation reached 99.2\%; it reached 68.7\% on N-CMAPSS and failed on bearings (0 of 17). In our data, useful alarms coincided with failure becoming visible consistently before the required notice and with monitored units that resembled those used for training. This is a post hoc reading of eight datasets, so both conditions, and the cost of the alarms, need to be checked for each fleet.

\section*{Declaration of generative AI and AI-assisted technologies in the manuscript preparation process}

\noindent During the preparation of this work the authors used Claude (Anthropic) in order to assist with code development, data analysis and drafting of the manuscript. After using this tool, the authors reviewed and edited the content as needed and take full responsibility for the content of the published article.

\section*{Data availability}

\noindent The PHM08, C-MAPSS and N-CMAPSS datasets are available from the NASA Prognostics Center of Excellence data repository, the battery data of Severson et al. \cite{severson2019} from \url{https://data.matr.io}, and the bearing data of the IEEE PHM 2012 challenge \cite{nectoux2012} from the FEMTO-ST institute. Code, hashed protocols, derived data summaries and the result files needed to reproduce every number are in a repository whose link is withheld for anonymised review.

\begin{thebibliography}{34}
\bibitem{lei2018} Lei Y, Li N, Guo L, Li N, Yan T, Lin J. Machinery health prognostics: A systematic review from data acquisition to RUL prediction. Mechanical Systems and Signal Processing 2018;104:799--834. \url{https://doi.org/10.1016/j.ymssp.2017.11.016}
\bibitem{si2011} Si X-S, Wang W, Hu C-H, Zhou D-H. Remaining useful life estimation -- A review on the statistical data driven approaches. European Journal of Operational Research 2011;213(1):1--14. \url{https://doi.org/10.1016/j.ejor.2010.11.018}
\bibitem{ramasso2014} Ramasso E, Saxena A. Performance benchmarking and analysis of prognostic methods for CMAPSS datasets. International Journal of Prognostics and Health Management 2014;5(2). \url{https://doi.org/10.36001/ijphm.2014.v5i2.2236}
\bibitem{depater2022} de Pater I, Reijns A, Mitici M. Alarm-based predictive maintenance scheduling for aircraft engines with imperfect remaining useful life prognostics. Reliability Engineering \& System Safety 2022;221:108341. \url{https://doi.org/10.1016/j.ress.2022.108341}
\bibitem{kamariotis2024} Kamariotis A, Tatsis K, Chatzi E, Goebel K, Straub D. A metric for assessing and optimizing data-driven prognostic algorithms for predictive maintenance. Reliability Engineering \& System Safety 2024;242:109723. \url{https://doi.org/10.1016/j.ress.2023.109723}
\bibitem{saxena2008} Saxena A, Goebel K, Simon D, Eklund N. Damage propagation modeling for aircraft engine run-to-failure simulation. In: 2008 International Conference on Prognostics and Health Management. IEEE; 2008, p. 1--9. \url{https://doi.org/10.1109/PHM.2008.4711414}
\bibitem{wang2008} Wang T, Yu J, Siegel D, Lee J. A similarity-based prognostics approach for remaining useful life estimation of engineered systems. In: 2008 International Conference on Prognostics and Health Management. IEEE; 2008, p. 1--6. \url{https://doi.org/10.1109/PHM.2008.4711421}
\bibitem{heimes2008} Heimes FO. Recurrent neural networks for remaining useful life estimation. In: 2008 International Conference on Prognostics and Health Management. IEEE; 2008, p. 1--6. \url{https://doi.org/10.1109/PHM.2008.4711422}
\bibitem{peel2008} Peel L. Data driven prognostics using a Kalman filter ensemble of neural network models. In: 2008 International Conference on Prognostics and Health Management. IEEE; 2008, p. 1--6. \url{https://doi.org/10.1109/PHM.2008.4711423}
\bibitem{zheng2017} Zheng S, Ristovski K, Farahat A, Gupta C. Long short-term memory network for remaining useful life estimation. In: 2017 IEEE International Conference on Prognostics and Health Management (ICPHM). IEEE; 2017, p. 88--95. \url{https://doi.org/10.1109/ICPHM.2017.7998311}
\bibitem{li2018} Li X, Ding Q, Sun J-Q. Remaining useful life estimation in prognostics using deep convolution neural networks. Reliability Engineering \& System Safety 2018;172:1--11. \url{https://doi.org/10.1016/j.ress.2017.11.021}
\bibitem{ariaschao2021} Arias Chao M, Kulkarni C, Goebel K, Fink O. Aircraft engine run-to-failure dataset under real flight conditions for prognostics and diagnostics. Data 2021;6(1):5. \url{https://doi.org/10.3390/data6010005}
\bibitem{grinsztajn2022} Grinsztajn L, Oyallon E, Varoquaux G. Why do tree-based models still outperform deep learning on typical tabular data? Advances in Neural Information Processing Systems 2022;35:507--520. \url{https://doi.org/10.52202/068431-0037}
\bibitem{friedman2001} Friedman JH. Greedy function approximation: A gradient boosting machine. The Annals of Statistics 2001;29(5):1189--1232. \url{https://doi.org/10.1214/aos/1013203451}
\bibitem{nectoux2012} Nectoux P, Gouriveau R, Medjaher K, Ramasso E, Chebel-Morello B, Zerhouni N, Varnier C. PRONOSTIA: An experimental platform for bearings accelerated degradation tests. In: IEEE International Conference on Prognostics and Health Management (PHM 2012). IEEE; 2012, p. 1--8. \url{https://hal.science/hal-00719503}
\bibitem{rengasamy2020} Rengasamy D, Rothwell BC, Figueredo GP. Asymmetric loss functions for deep learning early predictions of remaining useful life in aerospace gas turbine engines. In: 2020 International Joint Conference on Neural Networks (IJCNN). IEEE; 2020, p. 1--7. \url{https://doi.org/10.1109/IJCNN48605.2020.9207051}
\bibitem{severson2019} Severson KA, Attia PM, Jin N, Perkins N, Jiang B, Yang Z, Chen MH, Aykol M, Herring PK, Fraggedakis D, Bazant MZ, Harris SJ, Chueh WC, Braatz RD. Data-driven prediction of battery cycle life before capacity degradation. Nature Energy 2019;4(5):383--391. \url{https://doi.org/10.1038/s41560-019-0356-8}
\bibitem{saxena2010} Saxena A, Celaya J, Saha B, Saha S, Goebel K. Metrics for offline evaluation of prognostic performance. International Journal of Prognostics and Health Management 2010;1(1). \url{https://doi.org/10.36001/ijphm.2010.v1i1.1336}
\bibitem{barlow1960} Barlow R, Hunter L. Optimum preventive maintenance policies. Operations Research 1960;8(1):90--100. \url{https://doi.org/10.1287/opre.8.1.90}
\bibitem{dejonge2020} de Jonge B, Scarf PA. A review on maintenance optimization. European Journal of Operational Research 2020;285(3):805--824. \url{https://doi.org/10.1016/j.ejor.2019.09.047}
\bibitem{gebraeel2005} Gebraeel NZ, Lawley MA, Li R, Ryan JK. Residual-life distributions from component degradation signals: A Bayesian approach. IIE Transactions 2005;37(6):543--557. \url{https://doi.org/10.1080/07408170590929018}
\bibitem{nguyen2019} Nguyen KTP, Medjaher K. A new dynamic predictive maintenance framework using deep learning for failure prognostics. Reliability Engineering \& System Safety 2019;188:251--262. \url{https://doi.org/10.1016/j.ress.2019.03.018}
\bibitem{mitici2023} Mitici M, de Pater I, Barros A, Zeng Z. Dynamic predictive maintenance for multiple components using data-driven probabilistic RUL prognostics: The case of turbofan engines. Reliability Engineering \& System Safety 2023;234:109199. \url{https://doi.org/10.1016/j.ress.2023.109199}
\bibitem{lee2023} Lee J, Mitici M. Deep reinforcement learning for predictive aircraft maintenance using probabilistic remaining-useful-life prognostics. Reliability Engineering \& System Safety 2023;230:108908. \url{https://doi.org/10.1016/j.ress.2022.108908}
\bibitem{xie2025} Xie Z, Abdin A, Fang Y. Toward decision-oriented prognostics: An integrated estimate-optimize framework for predictive maintenance. arXiv preprint 2025;arXiv:2506.19698. \url{https://doi.org/10.48550/arXiv.2506.19698}
\bibitem{varma2006} Varma S, Simon R. Bias in error estimation when using cross-validation for model selection. BMC Bioinformatics 2006;7:91. \url{https://doi.org/10.1186/1471-2105-7-91}
\bibitem{kapoor2023} Kapoor S, Narayanan A. Leakage and the reproducibility crisis in machine-learning-based science. Patterns 2023;4(9):100804. \url{https://doi.org/10.1016/j.patter.2023.100804}
\bibitem{nosek2018} Nosek BA, Ebersole CR, DeHaven AC, Mellor DT. The preregistration revolution. Proceedings of the National Academy of Sciences 2018;115(11):2600--2606. \url{https://doi.org/10.1073/pnas.1708274114}
\bibitem{pedregosa2011} Pedregosa F, Varoquaux G, Gramfort A, Michel V, Thirion B, Grisel O, Blondel M, Prettenhofer P, Weiss R, Dubourg V, VanderPlas J, Passos A, Cournapeau D, Brucher M, Perrot M, Duchesnay E. Scikit-learn: Machine learning in Python. Journal of Machine Learning Research 2011;12:2825--2830.
\bibitem{clopper1934} Clopper CJ, Pearson ES. The use of confidence or fiducial limits illustrated in the case of the binomial. Biometrika 1934;26(4):404--413. \url{https://doi.org/10.1093/biomet/26.4.404}
\bibitem{mcnemar1947} McNemar Q. Note on the sampling error of the difference between correlated proportions or percentages. Psychometrika 1947;12(2):153--157. \url{https://doi.org/10.1007/BF02295996}
\bibitem{holm1979} Holm S. A simple sequentially rejective multiple test procedure. Scandinavian Journal of Statistics 1979;6(2):65--70. \url{https://www.jstor.org/stable/4615733}
\bibitem{page1954} Page ES. Continuous inspection schemes. Biometrika 1954;41(1--2):100--115. \url{https://doi.org/10.1093/biomet/41.1-2.100}
\bibitem{roberts1959} Roberts SW. Control chart tests based on geometric moving averages. Technometrics 1959;1(3):239--250. \url{https://doi.org/10.2307/1271439}
\end{thebibliography}
\end{document}
